\chapter{Combat}

\begin{multicols}{2}

\section{Turns and Actions}

\textbf{Start:} Combat occurs when a group decides to fight against another. A combat is played in rounds, which represent 3 seconds of game time each.

\textbf{Preparation:} At the start of combat, the miniatures are positioned in their locations, and each character starts with 6 AP by default. 

\skill{Cunning}{Awareness}{ Cunning is used to determine priority when competing to take a combat turn first. It is also used to gain tactical advantages in combat and assess danger.

\textbf{Turn order:} It is defined by a "first to ask, first to play" rule. If several players want to play at the same time, all interested make a cunning test, and the one who rolls the highest goes first. A character is forced to spend at least 1 AP and cannot take its turn twice in a row.

\textbf{ Feint :} is a cunning against cunning test between two characters in melee. It must be initiated during either character's turn. On a hit, the target is forced to take their turn and spend at least 3 AP on it. On a crit, they must spend at least 6 AP. On a miss or graze, the effect is reversed and the feinter must take their turn and spend the respective AP. Feint forces spending available AP, but cannot cause negative AP and cannot force an action surge.

\textbf{Analyze:} Cunning can also be used to assess an enemy. Spend 4 AP +1STA to Rest and to make a cunning check against an enemy to try to obtain relevant information about them from visible equipment or physical characteristics. On a critical success, learn important information for combat, dangerous items they carry, or special abilities they may have. On a success, obtain a vague reference to a dangerous item or the opponent's power level. Grazes only provide obvious information, and failures provide nothing.
}

\textbf{Surprise:} A combat can start with characters being surprised. In this case, all characters who are surprised cannot use their action surge on the first turn. A character is surprised if it does not detect its attacker before it begins attacking. 

\textbf{Ending a turn:} Once all desired actions have been performed or prepared, a character can end their turn. They can save as many AP as they want to use later in the same round.

\textbf{Action surge:} All characters are entitled to one action surge per round. There are 3 basic types of action surges: movement, combat, reaction and focus.

\begin{proplist}
    \item \textbf{Movement surge:} 3 STA to gain 6 AP that must be used for movement immediately and allow running until the end of the turn.
    \item \textbf{Combat surge:} 3 STA to gain 6 AP that can be used for attacks or movement and must be used immediately.
    \item \textbf{Reaction surge:} 3 STA to gain 6 AP, but can only use reactions until the end of the round.
    \item \textbf{Focus surge:} 1 STA to gain 2 AP. The AP is free to use for anything. This surge is required to use shooting weapons, casting spells and using professional abilities.
\end{proplist}

\textbf{End of the round:} A round ends when no one has AP left and the remaining agree to end the round. When the round ends, all characters get reset to 6 AP minus any negative AP they had. Any unspent AP is lost.

\textbf{Rest:} is an action that costs 4 AP and recovers STA by an amount equal to STA/4. The character is allowed to move 4 AP worth of careful movement while resting during their own turn. The AP loss from ending a round without spending all AP can be prevented by pooling AP from both rounds to Rest, which happens at the beginning of the new round.

\textbf{Environmental and ongoing effects:} Any effect that occurs periodically per turn or is environmental has its effect applied at the beginning of the round.

\textbf{Negative AP:} AP can become a negative number, but never as a result of performing an action or reaction. When the character recovers their AP, they must deduct the negative value from their maximum.

\textbf{Suffocation:} If the character cannot breathe at the beginning of a round, they lose 1 STA and cannot Rest. If STA is lower than -5, they drop unconscious. If it goes under -10, they die.

% \textbf{Delay the turn:} The player can spend 1 AP during their turn to delay it. This allows the character to regain their token and play again within the same round. AP regeneration only occurs once per round.

\textbf{Reactions:} These are actions that can be performed on another character's turn but must be triggered by something. An example is trying to evade an attack, which is triggered when being attacked. A reaction takes effect before the action that caused it.

\textbf{Preparing a reaction:} Any action can be prepared on a player's turn, but doing so costs 1 AP. This allows the player to create any trigger they can imagine for a reaction. They must have enough AP to perform the reaction when the trigger occurs. A prepared action can consist of a movement and a use of skill and can be repeated multiple times as long as the trigger happens again.

\textbf{Compulsions:} These are forced actions. They must occur immediately if there are enough AP. If not, they force action surges. If even that is not possible, the remaining AP are lost, and the compulsion ends.

\textbf{Severe compulsions:} Whenever a compulsion forces a character to perform an act that they perceive will cause physical harm to themselves or an ally, the compulsion is severe. In this case, the character has +5 to resist the compulsion.

\textbf{Adding more characters:} It is possible to add more characters to the combat after it has started. This can only be done at the turn of a round. If the added character is in stealth, make the appropriate stealth tests at the beginning of the round.

\section{Movement}

Movement actions spend AP in exchange for changing the character's position on the grid. A character must always end their turn in the center of a square or hexagon, which means that moving only a fraction of a space is impossible. It is possible to pass through a space occupied by another character as long as they do not resist the passage. 

Types of movement:
\begin{proplist}
\item \textbf{basic movement:} Simple walking movement.
\item \textbf{careful movement:} Slow movement used to reposition in battle or improve stealth and balance. 
\item \textbf{running:} Can only be initiated during a movement surge, but can continue running after accelerating as long as no turns of 90 degrees or more are made per running block. The character can run a partial amount smaller than its running speed, but AP cost is the same.
\item \textbf{getting up:} Removes the fallen condition.
\item \textbf{jumping:} A jump has a vertical distance equal to half the horizontal distance. If performed during a run, increase the base calculation of the horizontal distance of a long jump to AGI/3. Movement cannot be voluntarily interrupted in the middle of a jump. If performed backwards, the horizontal distance is halved.
\item \textbf{crawling:} is the only usable movement speed while \textbf{lying down}.
\item \textbf{swimming:} is the only usable movement speed while swimming. It also applies the \textbf{lying down} condition.

\item \textbf{trample:} The character attempts to pass through an occupied space. If the opponent does not move out of the way, they must perform a push action to stop the movement.
\item \textbf{mounted movement:} A character needs minimum training to be able to give commands to a mount. During mounted combat, the character's movement speed matches that of the mount. Running and a movement surge do not spend character STA. Mounting and dismounting costs AP equal to standing up.
\end{proplist}

\begin{stattable}{Movement Costs and Speeds}{@{}lcc@{}}
Movement & Cost & Speed \\
\midrule
Careful & AP & 0.5 m \\
\rowcolor{fallowtint}
Basic & 1 AP & 1 m \\
Run & 2AP & AGI/3m \\
\rowcolor{fallowtint}
Jump & 1AP+1STA & AGI/4m \\
Crawl & 1 AP & 0.33 m \\
\rowcolor{fallowtint}
Swim & 1 AP & 0.33 m \\
% Swim Fast & 1 AP & 0.5 m \\
Stand up & 5 -AGI/5 AP & 0 \\
\end{stattable}

\skill{Swimming}{AGI - 10}{ Represents the character's swimming ability. While in the water, make a swimming check against the DL of the body of water. On a critical or hit, swim perfectly. On a graze, lose 2PA+1 STA just to stay composed. On a failure, the character must use their movement surge just to stay above water or fall 1m.

The swimming check should not be redone more than once per round unless the swimming conditions change.

Characters attempting to swim with one hand receive a -3 penalty on the check.

A swimming check in a deep pool is considered a DL of -2. Swimming at a beach with waves is 0. Swimming in a swift river or choppy sea would be around 4. Swimming under a storm or in a raging river is around a 7.
}

\skill{Climbing}{AGI - 10}{ Make a climbing skill check to overcome a DL and climb a surface. Climbing checks can be made safely or recklessly.

The climbing speed is equal to the crawling speed but costs 1 AP + 1 STA. On a critical, it does not cost STA. On a success, they climb normally. On a near miss, decide to go twice as slow or give up. On a failure, the character simply falls from the surface.

A character who is climbing is considered lying down. An attack that interrupts the climber will require them to redo their climbing check with a -3 penalty. In the case of multiple climbing checks in a single turn, always take the worst result. Climbing with one hand incurs a -3 penalty.

To determine the difficulty of climbing, consider whether the surface is moving, the quality of the grips, and the distance to be climbed. An easy climb, with minimal chances of falling, is considered a DL of 5. A DL of 15 is reasonable for out-of-combat situations but quite challenging in a difficult scenario. DLs greater than 20 will require equipped and specialized characters and are unlikely to be accomplished during combat. As a general rule, consider that for every meter, the difficulty increases by 1.}

\skill{Balance}{AGI - 10}{ Skill to perform acrobatics, avoid falling, and mitigate damage from falling. Can be used actively with a skill check or passively when something tries to knock you down.

\textbf{Difficult terrain} Every time the character moves through difficult terrain, they must make a balance check. The DL is based on how slippery, unstable, long, and narrow the path is. The character decides the type of movement they wish to use, then makes the skill check. On a critical, all movement is inconsequential. On a success, only moving at a normal speed is safe, but not jumping or running. On a graze, moving at a careful speed is safe. On a miss, only crawling is allowed. If the character enters or moves in difficult terrain using a type of movement they did not manage well enough, they receive some consequence (usually falling).

\textbf{Fall skill test:} Falling requires a balance test to decide how well the landing or dive goes. On a critical on the balance test, reduce the damage by 50\% and prevent T0 and T1 injuries. On a hit, reduce damage by TGH. On a graze, takes full damage and falls prone, and on a miss takes full damage, falls prone and hits the head. Receiving a T2 injury or higher always makes the character prone.

\textbf{Choosing Fall DL:} When falling, the DL to fall properly is generally proportional to the fall height. Falls from small heights in a bad position are more difficult too. As a rule of thumb, DL is 1 per meter until a maximum of 25. As a common exception to the rule, DLs for 1m can be considered trivial at -5, then 0 for 2m. 

\textbf{Fall damage} For this damage estimation, consider that the landing spot is a flat, hard surface. To estimate fall damage, do 2x fall height x DM$^2$, Where the DM is the character's DM and the damage ignores armor. Maximum height is capped at 50m. Soft surfaces, like sand or water halve the damage.

Note: DM is squared because large creatures suffer more from falling. For size 2, instead of making 0.75$^2$, just use 0.5. For size 4, instead of 1.5$^2$, use 2.
}  

\section{Damage and Injuries}

\subsection{Injuries and Death}

\textbf{Injury level (IL):} Any time a character receives damage, they will increase their injury level. They will receive -1 as injury penalty for every multiple of their \textbf{injury tolerance (IT)}. The base value for injury tolerance is 10. The penalty affects skills shown in the afflictions section.

\textbf{Injury Capacity:} is the maximum amount of injuries a character can take without collapsing and going immobile. Its value is 4x IT.

\textbf{Death:} occurs when the character's injuries exceed 5x IT.

\textbf{Healing Injuries:} characters heal by resting well and using potions or medicine. The details are explained in the exploration chapter.

\textbf{Damage:} Any time a character receives damage, it must compare the amount received with its TGH for the damage type to determine the injuries caused. Armor increases the amount of damage that can be received without receiving wounds. The table shows how much IL they receive from injuries at each tier. 

\textbf{Damage Tiers:} 

\begin{stattable}{Damage Tiers}{@{}lccc@{}}
Name & Threshold & IL & Wound \\
\midrule
Tier 0 & armor & 1 & 0 \\
\rowcolor{fallowtint}
Tier 1 & armor+TGH & 5 & 0 \\
Tier 2 & armor+2xTGH & 10 & 50\% \\
\rowcolor{fallowtint}
Tier 3 & armor+3xTGH & 20 & 100\% \\
Tier 4 & armor+4xTGH & 30 & 100\% \\
\rowcolor{fallowtint}
Tier 5 & armor+5xTGH & 40 & 100\% \\
Tier 6 & armor+6xTGH & 50 & 100\% \\
\end{stattable}

\textbf{Interpreting IL:} A tier 0, or T0 injury means a small bruise or a patch of burned skin. A T1 injury is kin deep cut that bleeds a little or a large bruise. A T2 is a cut down to the muscle or a potentially cracked bone. A T3 is definitely a cracked bone and internal organ damage with bleeding. A T4 is near system shock, with multiple broken bones, impalation through the chest and potentially limb amputation. A T5 means instant death.

\textbf{Types of damage:} are: blunt, cutting, burning, electric, radiant, and corrosive. Objects and creatures may have resistance or vulnerability to each of these types of damage. Different types of damage cause different kinds of wounds. 

% \textbf{Multiple types:} Some attacks can cause multiple types of damage. In this case, the damage thresholds are calculated separately for each type of damage, and the injuries caused by each type are considered separately.

\textbf{Physical attacks:} have blunt and cutting damage components. If the weapon is not capable of cutting its target, the damage is blunt, otherwise, use the most advantageous of the two.

\textbf{Blunt damage:} is defended by armor protection. 

\textbf{Cutting damage:} is defended by armor RES.

\textbf{What cuts?} A weapon can cut something if their hardness is higher than their target's. There are 4 levels of material hardness. The first is for liquids. The second is for fabrics, flesh, and similar materials. The third is for wood, horns, bones and other solids. The fourth is for metals and rocks. More details on the gear section.

\textbf{Burn, radiant:} Their respective defense is INS. Dealing Tier 0 injury or higher leaves the target burning. 

\textbf{Corrosive:} Its respective defense is INS. Tiers 0 to II of damage leaves the target corroding at that tier of damage, which reaplies the damage until armor is removed. Higher tiers still apply Tier II of corroding.

\textbf{Electric damage:} It is also defended by INS. It causes interruptions when a character receives a tier 1 or higher damage from electric damage.

\textbf{Wounds:} can be caused by receiving high tier damage. Their healing counter is separate from the IL counter. 

\begin{proplist}
    \item \textbf{Broken arm:} Makes the arm useless. Equivalent to 10 IL to heal.
    \item \textbf{Broken leg:} Makes the leg lame. Equivalent to 15 IL to heal.
    \item \textbf{Bleed:} a special wound that requires healing to stop bleeding. For every STA spent and at the end of every round in combat, the character receives +1 to their IL per Bleed cure untreated. Bleeds caused by a T2 wound requires 1 heal to be healed. For every extra tier, one extra healing required. Bleeding is caused by damage to the chest or head. When combat/chase ends, bleeding causes 3x its current value and stops by itself when untreated.
    \item \textbf{Shocked:} cause immobile. Happens with damage to the chest. Has 50\% chance on a T4 hit and 100\% chance on a T5.
    \item \textbf{Unconscious:} causes unconsciousness. Happens when receiving damage in the head. Takes at least 1 minute to wake up without intervention.
\end{proplist}

\subsection{Additional effects:}

\textbf{Interruption:} is an effect that causes the victim to lose their action and associated costs when it occurs during an action. It causes a minimum loss of 2 AP when there is nothing to interrupt or the interrupted action costs 1 AP. It happens when a tier 1+ blunt or electric damage is inflicted or when the character is pushed. 

% \textbf{Knockback:} The victim is knocked 1 space away from the source of the attack. It is interrupted as well.

% \textbf{Knockdown:} The victim fall to the floor and becomes prone. It is interrupted as well.

\textbf{Stun:} An interrupt that has minimum loss of 4 AP.

\text{Success Overflow:} When a character receives a hit from an attack, they may receive additional effects, chosen by the attacker. Every point that the attack roll overflows the DL to hit posed by the defender during an attack can be used to activate these effects. These points are called Success Overflow Points(SOP). The effects are:

\begin{proplist}
    \item \textbf{Armor Bypass:} Only attacks from weapons with the precise property can use this effect. This costs a number of SOP defined by the targets armor deflection. It hits the inner layer of the armor, which has lower RES and is present in rigid armors. Penetrating is applied automatically if the attack also has the property.
    \item \textbf{Penetrating:} Costs a number of SOPs equal to the target's deflection to be able to use cutting damage against objects with the same hardness. For the same cost, ignores fiber armor RES entirely. This can only be done by weapons with the penetrating property. 
    \item \textbf{ Extra cut:} Increases cutting damage by 1xDM per SOP. Only weapons with the bladed property can use this.
    \item \textbf{Smash:} Costs a number of SOPs equal to the target's deflection to stun them if the damage is at least over T1 or interrupt if its over T0.
\end{proplist}
    
\textbf{Localized damage:} Changing location costs points based on the location. Each type of creature may have points that must be specifically aimed at when using the attack, called vulnerabilities. When someone declares that they are aiming at that specific point, the attack is modified according to the description of the vulnerability. 

To illustrate, here are some targets for humanoids. There are 4 locations: chest, hand, leg, and head.

\textbf{Chest:} Costs 0 points. Standard location. 

\textbf{Hand:} Costs 10 points, or only 5 when used against a target that blocks with a weapon. Hands have no armor unless the character is wearing gauntlets. A wound renders the hand useless. Maximum IL damage accounted from this is 5. Taking 20 IL damage amputates.

\textbf{Leg:} Costs 0 points. A wound leaves the character lame. Maximum IL damage from this is 10. Taking 30 IL amputates. 

\textbf{Head:} Costs 5 points. Getting a wound or stun to the head causes unconsciousness. The head has a 50\% chance of a wound on tier 1 and 100\% on tier 2 when receiving blunt damage. Tier 4 damage has 50\% chance of instant death, T5 has 100\%. Bypassing armor is possible with any weapon, unless the helmet is closed.

Some creatures may have their vulnerable point out of reach of others, requiring long weapons, jumps, and proper positioning to reach. These special cases can be added where the GM deems it makes sense. The ones added here are the most common and most tactically relevant ones.

\subsection{Afflictions}

Afflictions are the negative conditions that can affect creatures. 

There are 5 categories of affliction: Mobility, Injury, Vision, Mental, and Health. 

\textbf{Injury (IL):} penalties affect STR, AGI, STA. This does not affect movement speeds, nor TGH. It affects all other usages of attributes.

\textbf{Sensory:} Strike, Defend, Accuracy, Reflex, Detection, Balance, Climb, Exploration, Stealth and Prestidigitation.

\textbf{Mental:} Penalties affect Exploration, Cunning, Will, Insight and all Knowledges and all spellcasting.  

The list of afflictions is:

\begin{proplist}
\item \textbf{Lame:} Cannot run, jump, or perform basic movement.
\item \textbf{Prone:} Can only crawl and stand up. Takes -3 to any attempt to evade.
\item \textbf{Grappled:} Movement requires pushing. Focus is locked to the main grappler.
\item \textbf{Immobile:} Cannot move and cannot use any combat or movement skills other than escape. The SD is reduced to -5 + SMx1. 
\item \textbf{Disoriented:} Gives -2 to sensory penalties.
\item \textbf{Oblivious:} Gives -5 to sensory penalties.
\item \textbf{Blind:} Cannot use vision.
\item \textbf{Deaf:} Cannot use hearing. 
% Requires a balance check for basic movement on any terrain.

\item \textbf{Burning:} Burning damage ignores armor insulation. Ends automatically at the beginning of the next round except when reapplied. 
\item \textbf{Corroding 0/I/II} Applies a tier of injury every round. Can only be ended by specific items, jumping on water or removing the armor. Cannot be applied to unarmored.

\item \textbf{Unconscious:} The character cannot do anything other than Rest and Will and Health tests. They can also make detection and sensing tests, but receive -10 on these tests. These tests can be made to wake up. The DL and the allowed test depend on the situation. Unconscious characters breathe automatically and cannot hold their breath. 

\item \textbf{Afraid:} Receives -1 Mental and cannot make a combat surge. The fear ends when the character succeeds on their Will test or becomes angry. 
\item \textbf{Enraged:} Receives -1 Mental and cannot make a reaction surge. The anger ends when the character succeeds on their Will test or becomes afraid. 
% \item \textbf{Charmed:} Receives -1 Mental and behaves like permanently distracted. Cannot coexist with afraid or enraged.
% \item \textbf{Distracted:} Receives -3 to Detection and -3 Cunning.
\item \textbf{Dominated:} Under the control of another character. Must obey commands for the duration but can make a Will test to break out of domination whenever a command constitutes a severe compulsion. Outside of combat, can make a test once per hour.

\item \textbf{Intoxicated I/II/III} -1/2/3 Mental.
\item \textbf{Tired/Exhausted/Confused} Penalty of -1 mental/-2 mental/-4 mental. Becomes vulnerable to domination attacks. It is possible to be confused without being tired or exhausted, but extreme exhaustion causes confusion.
\item \textbf{Weakened/Malnourished} -1/-2 Health. -1/2 to IT.
\item \textbf{Thirsty/Dehydrated} -1/-2 Health, -1/2 to IT
\item \textbf{Sick} -2 Health. -2 to IT.
\end{proplist}

\skill{Health}{ CON }{ Represents how strong the character's body is against diseases, poisons, exposure to the elements, and how quickly they recover from injuries and exhaustion. 

The survival chapter describes how health works for healing.
}

\section{Melee Combat}

%---------------------ATTACK

\skill{Strike}{Melee}{ Striking consists of a melee weapon attack on a target using the strike skill test against the opponent's defense. The AP cost is 3 AP, but that may be modified by attack variations. Grazes deal 50\% damage and misses do nothing. Hits deal full damage, and any extra points over a hit can be used to cause additional effects.

\textbf{Attack against equipment:} An attack can be made against a weapon, shield, or any object in order to break them. Objects break when their RES Rules on equipement breakage can be found in the equipment chapter.

Other variations of melee attacks are available depending on the type of weapon used. The type of attack must be declared when used. Each attack type must be used individually.

\textbf{Heavy Attack:} It is possible to make a heavy attack with weapons that have the "heavy" property. Heavy I increases damage of physical attacks by +0.5xSTR, but have an added cost of +1AP. Heavy II adds +1xSTR, costs +2AP+1STA, and gives a -2 penalty to hit. Heavy III adds +1.5xSTR, costs +3AP+1STA, and gives a -3 penalty to hit. Bonus applies to blunt damage and cutting damage of piercing or bladed weapons.

\textbf{Sweeping Attack:} spends 1 additional AP on the attack and hits everything in a semicircle. Decide whether to attack from right to left or vice versa and attack the enemies in that order in a 180-degree arc. If the attack is blocked or redirected, the impact is reduced for all. If one character is behind another, only the closest is hit. Sweeping attacks do not hit fallen enemies. There is no size bonus or penalty to strike in sweeping attacks.

\textbf{Opportunity Attack:} A melee weapon attack that can be used as a reaction when an opponent attempts to perform a triggering action. It is possible to cancel the triggering action and reuse the AP spent to defend against an opportunity attack, but not with evasive jump. The attack requires the normal AP cost and has a +2 bonus to hit. Triggering actions include casting spells that cost 3 AP or more, using ranged weapons, using a standard action, standing up in melee range, and moving towards a melee weapon while within its attack range. The attack occurs before the effect of the triggering action, unless stated otherwise.

\textbf{Braced Attack:} A melee attack that attempts to increase strength by firmly bracing the weapon. It can be used as a reaction when a target is moving towards the weapon or as an action in situations where the attacker is mounted and running. The attack happens when movement is between two spaces within weapon range. It costs an additional +2AP and +1STA on top of the normal attack cost and has damage bonus of +1.5x STR on a hit. If the damage is blunt, this also pushes the target.

\textbf{Hook Attack:} A melee attack with a hook weapon that attempts to knock down the opponent. Gains +STR/2 damage if the opponent attempts to perform an evasive jump and +STR if they are running. This can be used as a reaction against running targets that move away from the weapon within two spaces, which are both inside its melee range. If the damage is blunt, this counts as a drag, from grappling.

\textbf{ Assassinate :} Requires being in short range and having a free hand or equipping a small weapon on the off hand. This costs 1 extra AP on the normal cost of the attack. Precise weapons gain the +3 to bypass armor. This is equivalent to opening a gap in the oponent's armor with one hand and attacking with the other.
}


%-----------------------------------DEFEND

\skill{Defend}{Melee}{ This is the defense against the character's melee attacks. Whenever attacked, the character must spend AP to use the defend with some object or part of the body; otherwise, they use the SD. If the strike hits or crits, damage is taken normally. If it grazes or misses, the result depends on the type of defense chosen. There are four types of defense: Evade, Evasive Jump, Intercept, and Block.

% \textbf{Flinch:} This is a reaction that costs 1 AP and increases deflection by 2. The attack is still made against the SD.

\textbf{Evade:} The defender spends 2 AP to try to get out of the way of an attack. On a graze, the attack has 50\% damage, and on misses, it has 0.

\textbf{Evasive Jump:} is a modification on evade. It costs an extra 1 STA and allows the character to jump from the attack while evading an attack. This movement cannot be followed and gives +AGI/3 on the skill test. This can only be used if there is space to jump. The jump can be backwards or sideways and uses the movement speed of jumping backwards.

\textbf{Block:} attempts to block an attack by spending 2 AP and using any weapon attack or shield that contains the DEF property. On graze, the attack damage is reduced by the block value. On a miss, by 1.5x as much. A blocking object can break if damage is greater than RES. 

\textbf{Intercept:} uses 3 AP to intercept the blow. It requires the defender to be in short range and can be done with anything that has the DEF property. On a graze, stops attacks unless the attacker has 5 grapple higher than the defender. On a miss, stop it if the attacker has up to 8 grapple more than the defender.

\textbf{Protect:} It is possible to defend an ally by staying within one space distance of the line between attacker and target. The defender can block or intercept an attack directed at an ally.}

%----------------------- Grapple

\skill{Grapple}{Melee + STR-10}{ Used to grab and subdue others. When a character is under the effect of grapple, it cannot move without dragging the other grappler. Additionally, the list of things a grabbed character can do is limited to the points below:

\textbf{Initiate the Grab:} attackers must make a strike test with a weapon that has the grapple property against the opponent's defense or SD. On a hit or critical, the opponent is grappled, and any movement initiated by them is stopped. It is possible to grapple back automatically just by having a weapon with grappling property equipped, like a free hand ( a two handed weapon can be held with one hand ). If a creature's grapple skill is 10 points higher than its oponents, it does not suffer the grappled affliction.

\textbf{Attack and Defend:} Strikes between opponents involved in a grapple can be done with short range attacks. It is not possible to evade or block attacks, only intercept. Once grappled, a character will center its focus on its grappler and cannot shift it to another for the purposes of flanking o blocking projectiles. Any weapons with the grappling property can use the grapple skill instead of strike to attack during a grapple.

\textbf{Grapple Maneuvers:} These are performed during a grapple by any of the participants in order to \textbf{escape, move, immobilize, disarm, or knock down} an opponent. These actions cost 3 AP +1STA, must be defended with the grapple skill, and require the defender to spend 2 AP+1 STA or suffer a -5 penalty. If the grapple attack has any damage, it deals that damage whenever a grapple maneuver is used against the opponent.

\textbf{Drag and Push:} On a hit, moves the opponent one space x MM. A crit in the grapple check makes the opponent prone. Pushing can be used simultaneously, without exta AP cost and sharing the dice roll, with initiating the grapple if the strike test is a hit. It gets +3 to the grapple test if the grapple initiated while running or when the hooked or braced attacks get STR bonus.

\textbf{Escape:} Escapes from the grapple on criticals and hits. Being interrupted allows for a reaction to escape without the possibility of active resistance.

\textbf{Disarm:} Removes something from the opponent's hands on a critical. It is possible to grab the opponent's weapon on a hit, preventing them from using it until they manage to escape.

\textbf{Knockdown:} The opponent falls to the ground on a critical. It is possible to throw oneself along to achieve a knockdown on a hit.

\textbf{Immobilize:} The opponent becomes immobilized on a critical. It is possible to stay immobilized yourself to achieve immobilization on a hit.

Multiple characters can assist each other in the grab. Any skill test made by several characters uses the highest skill value among them plus 3/2/1 for each additional character, up to a maximum of 5. All characters involved must spend the necessary AP. Characters of a smaller size category add a maximum of 2.
}

%----------------------SHOOT

\section{Ranged Combat}

\skill{Accuracy}{ Ranged Combat }{ Accuracy consists of a throw or shot directed at a target using the Accuracy skill test in ranged combat against the opponent's reflexes or their SD, should they choose not to react to the attack. The cost in AP depends on the weapon used and the type of attack performed. Grazes deal 50\% damage and misses do nothing. Hits deal full damage, and any extra points over a hit can be used to cause additional effects.

\textbf{Distance Measurement Recommendation:} Measuring distance is a time-consuming process, so it is more practical to estimate the distance intuitively and give the benefit of the doubt in favor of the shooter.

\textbf{Vertical distance:} The vertical range of a ranged weapon is half the horizontal range.

\textbf{Canceling Attack:} It is possible to evade, dodge, or block while shooting, but the shot is canceled, and the costs of AP and STA are the greater of the defensive action or the shot.

\textbf{Immobile Targets and Shooting Errors:} The DL to hit an immobile target occupying a space of 1m is equal to -5-1*SM. If a projectile misses the test, it hits some adjacent space or a distance even greater for every 3 points the test is below the DL.

\textbf{Moving Targets:} It is possible to shoot at moving targets as a reaction. The total amount of AP spent to shoot cannot exceed the amount of AP spent on movement. The AP count starts when the target becomes visible. The target's position at the moment of the shot is where it is at the end of the movement.

\textbf{Prepared Shots:} If a character prepares a ranged attack, the shot occurs when the target becomes visible. It is possible to prepare to shoot the same target as an ally as a reaction to their ranged attack. The shot occurs at the same time as the ally's shot.

\textbf{Coordinated Shots:} Whenever an ally shoots a target, it is possible to target the same target in the same action, as long as they are in range. The target defends all shots with a single action.

\textbf{\large Types of Ranged Attacks:}

Other types of ranged attacks will be available depending on the type of weapon used. The type of attack must be declared when used, and only one can be used at a time unless stated otherwise.

\textbf{Throw:} Basic ranged attack with a throwable weapon. Throwable weapons gain range based on STR.

\textbf{Shoot:} Basic ranged attack with a shooting weapon. Requires Focus surge to use. Shoot is only allowed against targets within 30m.

\textbf{Snipe:} Modifies Shoot to cost an additional 2 AP, but allows shooting targets at any distance.

\textbf{Quick Shot:} Modifies Shoot to cost 1 AP less, but can only be used within 10m of the target. 

\textbf{Ranged attacks outside the combat area:} Any ranged attacks with range higher than 50 can hit outside the combat area at a penalty of -5.
}

\textbf{\large Explosions:} This type of attack results from something like a bomb, which radiates from a central point and loses strength with distance. Explosions are defended against with a reflex test. If the explosion comes from a projectile, the DL of the explosion is equal to the shooting skill. If the explosion occurs before being perceived, no test can be made. If the bomb is perceived, players receive a certain amount of AP to move away before it occurs.

Explosions have different magnitudes of effects depending on where things are in relation to the center. The center is where the axis of the spell's radius is located. Anyone caught in the center of the radius receives the critical effect, while those in the middle are hit, and the ones in the edge are grazed. If a creature occupies multiple spaces, apply the strongest effect. For very large explosions (>5m radius), the critical and graze zones can be expanded (30\%/40\%/30\%). 

The damage from explosions is 200\% on a critical, 100\% on a hit, 50\% on a graze, and 0 on a miss. Effects of an explosion that have a DL, such as pushing enemies, receive +5 against enemies in the critical area and -5 against those in the graze area.

\textbf{\large Sprays:} They are a type of explosion, with the same damage progression based on distance from the origin and the same DL definition mechanic. Sprays target all characters in range, which their reflex saves to escape, and move their characters according to their result. The attacker can choose the exact direction of the cone after the movement.

\textbf{\large Huge projectiles:} Draw the projectile's trajectory and the places it occupies until stopped. This is defended by reflexes. 

\textbf{\large Objects as explosions:} A large object can be considered an explosion. Draw its trajectory and the places it occupied until it stopped. This is defended b reflexes, but only the outer edge counts as a graze. 

%--------------------- REFLEX

\skill{Reflex}{ Ranged Combat }{ Used to react as quickly as possible to something. This can be used to avoid explosions, projectiles, and other effects that require quick reactions.

\textbf{Evasion:} When the character is targeted by projectiles, they can spend 2 AP to react. This AP cost is used to move and does not need to be spent again, but any STA cost must be paid. Upon receiving a critical or hit, take the attack normally and then use up to 2 AP to move if not interrupted. On a graze, the character can jump to try to gain cover before being hit and only takes half damage if they fail. On a miss, the character does not take the attack and can spend their movement surge immediately to escape or do the same as in the graze.

\textbf{Guard:} If using a shield, they can spend 2 AP to block ranged attacks against themselves or adjacent characters, as long as they are closer to the projectile source than the adjacent character. Shield deflection is added to guard as a bonus. The damage is reduced by the shield RES on grazes and misses. Hits and crits damage the character directly.

\textbf{Cover:} When using a large shield, it is possible to take total cover behind the shield after using guard. The cover is canceled if the character performs any action other than careful movement. This cover only extends to allies who stand in the same space.

\textbf{Avoiding an Explosion:} Spend 3 AP and make a reflex skill test against the DL of the explosion. On a critical, the character can run by spending one extra STA. On a hit, they can spend an extra STA to jump in any direction before the explosion occurs. On a graze or miss, they can move 2 AP after the explosion.

\textbf{Explosion Block:} Can only be done with a shield. Make an active reflex + shield deflection test against the DL of the explosion. If successful or critical, the explosion damage gets reduced by the defenses from the shield. The impact from explosions is unblockable. For explosions that deal cutting damage, the damage is not doubled and a hit blocks half of the damage instances and crits block all of them.

In all cases, if the character is interrupted, they cannot move after the interruption.
}

\section{Positioning and Visibility }

\textbf{Flee:} It is possible to activate the movement surge as a reaction for a character to prevent them from entering melee range, getting any closer in melee or after receiving a melee attack. This reaction interrupts the oponents turn.

\textbf{Follow:} A character can, as a reaction to any movement except running, follow another character who is already within melee range. The action must be performed simultaneously with the triggering movement and cannot cost more AP than it.

\textbf{Flanking:} Whenever two characters cannot fit completely within a radius of 180 degrees relative to a target, the target must focus on one target at a time. Use 3 surrounding hexes for a hex grid. Switching focus costs 1 AP and is necessary in order to react to the target. If the target of the attack moves as a part of its defense, like with evasive jump, the 1 AP cost is waived as long as the final destination reduces the angle back to 180 degrees or lower.

\textbf{High Ground:} This happens when one character is standing on terrain around 1m higher than its surroundings. One character no longer offers cover against ranged attacks to another if one is in front of the other against shots from the high ground. The melee reach is reduced by 1m unless the person on the lower ground targets the legs. The one on the high ground has the reach reduced to hit the legs of someone on the lower ground. Both receive a +2 bonus to their melee defense against each other. A difference in elevation is a defensive tactical resource where whoever has the high ground has a slight advantage.

\textbf{Cover:} A character can protect themselves against projectiles and explosion attacks by moving behind something, such as an ally or a wall. Staying behind something acts as a shield, reducing the damage taken by a certain amount. A character has cover if the shortest path from the origin of the explosion or projectile to a destination passes through a space containing a blocking object. Cover blocks vision unless it is transparent. Cover can be partial. In that case, it does not block vision and provides a bonus, usually +3, for reflexes and absorbs damage as if it were a shield in the case of grazes and misses. 

\textbf{Visibility:} is a temporary characteristic of a target. Can be good, bad, or zero. Bad visibility applies -2 visibility penalties to anyone targeting them, as if they were disoriented. Zero visibility means that the target is untargetable through vision. The terrain can offer visibility status to those in it through fog, smoke and dim light.

\textbf{Darkness and light:} In a situation of darkness, the field of vision is determined by the light source. If light source is at a distance, consider that the field around it is visible. Anyone in the light radius is in good visibility and anyone out is considered at zero visibility. 

\textbf{Fog and smoke:} Opaque gases reduce vision radius of anyone inside of it to 2m by default and makes them zero visibility to those outside. Standing at the edge of the area makes them bad visibility to those outside, but allows good visibility for the one inside.

\textbf{Underwater Combat:} Any attack that does not have the precise tag or the grapple tag has its damage reduced by half. Fire, radiant and acid damages are also reduced by 50\%. The range of projectiles is reduced to 3m for small objects (the size of arrows), 5m for medium ones (the size of spears) and 8m for large ones (the size of ballista shots, since larger projectiles decelerate less). All explosions have half the radius, but blunt damage is doubled.

\textbf{Mounted Combat:} Attacking in melee receives a mount-dependent penalty while mounted, usually around -2. Accuracy receives the penalty while the mount is moving, not when standing still. Any movement, except careful movement, is considered running to trigger reactions. If the character is forcibly displaced while mounted, they will be dismounted. The only defenses against melee attacks allowed during mounted combat are block and redirect.

\section{Interactions with objects}

This section is a guideline for interactions with objects in combat. The exact values can be modified by the GM. This is just a basis to work with.

\textbf{Lifting Weight:} The basic rule of strength is that you become 3.33x stronger for every 5 points of strength, following the VM from the size categories. The exact weight lifted will depend greatly on the exact lifting movement. It is fair to consider that a regular human with strength 10, for any random object shape, can lift about 20kg within a 3 AP + 1STA action. That is about one size category lower than its own.

\textbf{Pushing objects:} To push an object, or when struck by one, use the push action as described in Grapple, but use STR as a skill. A hit stops the object and moves it in any direction by 1m per 3 SOP. When struck by an object, in a graze, the object can be redirected by up to 60 degrees.

The push DL depends on weight. When stopping an object, also consider its speed. To help decide on a DL, consider that 50kg (size 3) is the base dificulty, with DL 0. Change the DL by 5 for every size category. Increase the DL for higher speeds. A suggestion would be +1 per m/s for objects of size 3. Reduce it for smaller objects and increase it for larger ones.

\textbf{Defending against objects:} Small objects are considered projectiles and defended as such. Large objects can be considered explosions on a line given by their path. Use reflex to dodge them. Consider how difficult it is to dodge subjectively, using the basic guidelines on how to choose a DL.

\textbf{Damage from Impact with objects:} Depends on what is hitting what, at what speed, and whether the target is fixated or free to move. \textbf{Action and reaction:} The impactor takes the same damage as the target. In practice, it is recommended that the value should be estimated quickly and intuitively at gametime. Use the instructions ahead to whatever extent it is useful.

\textbf{What objects to consider:} Objects with speeds significantly higher than 20m/s and sharp objects in general are considered weapons and should have their damage estimated separately.

\textbf{The base damage estimation:} The initial assumption for these numbers is that the object is rigid and blunt. The damage is 10 x DM x speed category, where speed categories are 1 for 5m/s, 2 for 10m/s and 3 for 20m/s and DM refers to the impactor's size category DM. This is also used for damage between objects when the target is fixated in place, like a door or wall.

Adjust the damage as follows:

\begin{proplist}
    \item \textbf{Freedom of movement:} If the target is free to move, it is swept along instead of crushed: halve the damage. Do not halve the damage for objects 2 size categories below the target.
    \item \textbf{Contact:} soft or wide objects (bodies, sand) halves damage.
    \item \textbf{Armor:} Damage from objects of size equal or greater than the character ignore armor. Note: armor spreads damage, but large objects have enough force to damage even when energy is spread out.
\end{proplist}

\skill{Prestidigitation}{ DEX }{ Represents a character's manual dexterity. It can be used to attempt to steal an object without being noticed, pick a lock, write neatly, conceal small objects, perform delicate operations with the hands, and other manual skills.

\textbf{Standard Action:} Consider that 1 AP is equal to half a second and estimate how long an action should take. In general, taking an item from a belt or strap around the armor should take about 3 AP, while opening a backpack and taking something costs about 5 seconds. The AP cost of a standard action is 3 AP.

\textbf{Draw/Put away weapons:} Weapons with the draw property take 1 standard action to be drawn. The rest take 2.

\textbf{Conceal Objects:} Used as a defense against detection tests to identify objects carried by the character. This is a detection vs prestidigitation test and works like a regular detection test. There is a -5 penalty to conceal items for every item size larger than small.

\textbf{Theft:} A standard action can be used to steal an item from someone's body or place an item on them. This is done by having one hand free and using a standard action to perform a Prestidigitation skill test against the opponent's Detection +5x for every item size category over small (relative to who must detect it). On criticals, it is stolen without being noticed. On successes, the character can choose between successfully stealing or failing without drawing attention. On grazes and failures, attention is drawn and it fails. Stealing an item from stealth gives a bonus of +5.

\textbf{Pick Lock:} These tests require knowledge of mechanics and receive -5 for each level below the requirement. On criticals, it takes 5 AP. On successes, between 10 seconds and 1 minute. On a graze, it takes between 10 minutes and 1 hour, and on a failure, it is impossible to open.

When Prestidigitation reaches 5, the cost of the standard action becomes 2. If it reaches 10, it becomes 1 AP. If Prestidigitation is less than -5, they cost 4 AP.
}

\section{Environmental Hazards}

These are effects that play a role in controlling terrain, forcing actions out of oponents and potentially causing damage.

This list shows the possible effects and how to handle them. The actual item list is in the gear section.

\subsection{Fire}

Fire is one of the most available ways to force characters to move. It is not a very fast source of damage, but it cannot be ignored.

Fire causes burn damage at the beginning of every round to anyone in it as an environmental effect. Any creature that occupies more than 1 burning space takes the damage multiplied by the amount of burning spaces it is in. 

A burning surface lasts at minimum until the end of combat, or until extinguished.

Here are a few common flammable surfaces and the amount of burning damage they cause when covering the same space as a character. 

\begin{stattable}{Flammable Surfaces}{@{}lcc@{}}
Fuel & Burn/round & Spread \\
\midrule
Dry grass, thatch & 5 & 1/round \\
\rowcolor{fallowtint}
Burning cloth & 5 & 1/round \\
Burning Oil/Alcohol & 8 & 1/round \\
\rowcolor{fallowtint}
Wood pyre, timber & 8 & 0 \\
\end{stattable}

\textbf{Extinguishing:} Fires can be extinguished by muffling it, throwing water or with magic. Both muffling and throwing water require 4 AP by default and being in a burning space or adjacent to it. Not every fire is extinguishable through water or muffling. Some fires require more than 4 AP to be extinguished.

Optional: \textbf{Fire spreads} each round to adjacent spaces that contain fuel. Adjust spread speed depending on fuel.

\textbf{Sticky fuel:} lasts until extinguished and follows mobile targets. Extinguishes automatically at the end of combat unless otherwise stated.

\subsection{Gas}

Gases can have several effects: Poison, explosion, suffocation, reduced visibility and blindness. It lasts until the end of combat or dispersed somehow through wind. Muffling or extinguishing the source of the gas causes it to disappear in the beginning of the next round. Doing that costs 4 AP. Relocating the source moves the area somewhere else starting next round.  

\begin{proplist}
    \item \textbf{Suffocation:} anyone that starts the round inside a suffocating gas is suffocating by default. Abilities and items may allow breathing.
    \item \textbf{Reduced visibility:} Fog and smoke disrupt visibility as explained in Positioning and Visibility section.
    \item \textbf{Inhaled Poison:} Causes a poison effect to anyone who begins the round inside the gas or takes a Rest action. If the poison gas is noticed, it is possible to self inflict suffocation to avoid it. In order to detect poison gas, use an appropriate knowledge test to recognize that it is poison and/or a detection test to notice its smell if the smell is weak.
    \item \textbf{Blinding Gas:} Staying inside it for any amount of time calls for a health test to avoid blindness. Use a reflex test to avoid it. Protective gear over the eyes also prevents blindness.
    \item \textbf{Explosion:} Some gases explode in contact with fire or electricity. The damage depends on the gas and uses the rules for explosions.
\end{proplist}


\textbf{White Smoke:} affects visibility, but only suffocates in enclosed places filled by it. \textbf{Black Smoke:} Affects visibility and also cause suffocation.

\textbf{Smoke Spread:} smoke only spreads indoors. The spread is counted in minutes as how long it takes to fill a 10m by 10m room.

\textbf{Wind:} Requires balance or strength test to avoid being pushed. Clears the other gases for the round. If the gas has an active source, it comes back once the wind stops.

\subsection{Liquids}

\textbf{Liquid Spread and splash:} Smalls amounts of liquid (under 50 L) are considered to spread instantly and do so with the geometry of an explosion. Liquids spread around solids, not through them. Splashing a bucket of a liquid counts as a thrown ranged attack with 3m x MM range.  

\textbf{Corrosion:} Being soaked in corrosive liquids cause a corrosive damage. The damage is multiplied by the amount of corrosive filled spaces the character is in.

\textbf{Waves and Currents:} Requires strength test to avoid being pushed. If the character is swimming, they get pushed automatically, but can spend AP to swim and cancel the movement.


\subsection{Land and collapse}

\textbf{Slippery Surfaces, swampy terrain:} trigger balance tests. Hint at the difficulty of the test if the character noticed that the terrain was difficult.

\textbf{Traps:} Require detection tests. Gets Building or Survival knowledge bonus depending on where the trap is. Magical traps may use the appropriate magical knowledge.

\textbf{Collapse:} Works like an object as an explosion. Gets buried in the rubble unless the explosion grazes or misses. Escaping requires a Strength test.

\textbf{Jagged terrain:} causes crushing or piercing damage per movement. Balance test lowers the amount of damage.


\section{Chases}

Chase scenes happen when two or more teams are at risk of entering combat. It can happen as both teams approach each other or as they are chasing each other.

\textbf{Escape:} The chase starts when everyone in a party decides to run away. It always starts at the beginning of a new round. Any character that cannot run for whatever reason at the beginning of the round will prevent the chase from beginning, unless they get left behind. If they do, finish the fight, and then start the chase.

\textbf{Approach:} This starts when characters have range to attack the other team. The weapon range will decide at what distance both teams start from each other.

There are four distance tiers: Combat, Shooting, Far and Horizon.

\begin{proplist}
    \item \textbf{Combat:} This is under 50m distance. Can be smaller when indoors (a room) or in places with many obstacles.
    \item \textbf{Shooting:} This is from around 50m to 200m distance. This range can be smaller in a city, a dense forest or any other place with lots of hiding spots and obstacles.
    \item \textbf{Far:} This is from around 200m to 500m. This is the limit of what can be seen in many places. It will be smaller in places with many obstacles and hiding spots.
    \item \textbf{Horizon:} This is from 500m and beyond. This distance is only used in open terrain.
\end{proplist}

\textbf{Chase actions:} Chase actions can be running and scouting, ranged combat or hiding. 

\textbf{Escaping:} There are three ways to escape:

\begin{proplist}
    \item \textbf{Outrunning:} If the distance is already far or Horizon and the winner decides to increase the distance, the chase is over.
    \item \textbf{Hiding:} If the distance is at least far and there is a stealth event, in which the quarry succeeds, it escapes.
    \item \textbf{Obstacle:} If there is an obstacle that the quarry traverses and that the pursuers cannot or do not want to traverse.
\end{proplist}

\subsection{Running}

Running is how both teams get closer or farther from each other. It works in rounds, where both teams compare their speeds, and the fastest decides to close the distance or move away.

\textbf{Running at combat range:} This is a dash. Both sides use their movement surge to sprint, spending STA to do it. Whoever has the highest running speed wins. A draw does not change distance tier, but spends STA from the movement surge.

\textbf{Running at longer ranges:} This is a sustained pace running. The running speed is STA + AGI. Whoever is higher wins. 50\% chance on a draw. If the character is bleeding, it bleeds as if they have used an action surge.

\textbf{Catching:} The quarry is caught when one of its members does not have STA to run, or when it ends a round with any mobility affliction while in combat range. In this case, a combat or social situation starts. When the combat starts, the tokens can be arranged in any way, as long as the distance from the character who ended the chase is within 10m of one enemy and within 30m of any member of their team. If the catch happened over a grapple, the characters involved start adjacent to each other.

\textbf{ Routing :} When the pursuers are in combat range, but the quarry is not caught yet, everyone gets 6AP, but the pursuers have +2 bonus to all uses of strike,  accuracy and cunning. This works like normal combat, but no grid is pulled yet. Everyone is considered to be in range to attack or defend anyone. Once the AP are spent, the round is over and a new running test is made.

\textbf{Group chase:} In a group chase, use the speeds of the slowest in each group as the speed for the other team to beat.

\textbf{Separating the group:} If one participant in the group is too slow, it is possible to leave them behind. From now on, the distant members are one distance tier away from their group. If combat starts, the distant group takes 2 full rounds to arrive, and when they do, they do it at 30m distance from a chosen ally.

\subsection{Skirmish}

This stage can only happen at Shooting distance.

One team may decide to fight to attempt to injure the other before they get closer or to snipe enemies early. There are two possible situations: Both teams shoot at each other or one team shoots and the other attempts to close in. Either way, any weapon with a range higher than 50m can be used. Everyone is targetable by anyone, except when one whole team is invisible to the other.

\textbf{Both teams shoot each other.} This happens until one team runs out of ammo, someone changes strategy and decides to run, or a whole team dies. Everyone has 12 AP to spend as they like per round.

\textbf{One team attempts to approach.} The approaching team gets 6 AP and must spend a movement surge for running. The shooting team gets AP depending on the running speed of the slowest character of the approaching team. They get 18 AP for running speed 2m per 2 AP, 12 AP for 3m, and 9 AP for 4m or higher (Optionally, $36/speed$ AP). At the end, both teams will be in combat range. The approaching team can split and leave someone behind to shoot, in which case they get the same amount of AP as the attackers. 

\textbf{Split team combat:} If two teams are engaging in combat range, anyone at shooting range can still shoot from far away. They have their turn normally and are considered to be outside the combat area for the ranged penalties. If someone wants to go after them, they must end the round 30m away from any enemies (or their corpses). The combat just gets two areas within the same round order, where the new combat area starts with the two groups 30m apart.

Ranged combat can happen over terrain features, benefiting either team.

\textbf{Cover:} If a character has cover, they can decide not to leave it and stay untargetable, but then they cannot target any enemy with attacks or spells. For the ones who decide to shoot or move, cover is partial, with a bonus of +3 to reflexes by default.

\textbf{Visibility:} If visibility is affected by fog or darkness, use bad visibility if at any point the opponent becomes targetable. In case they never become visible, they stay at zero visibility. 

\subsection{Hiding and Terrain}

\textbf{Scouting:} The fugitives can look for a place to hide or a vantage point for combat using their detection skill. Only the fugitives can do this, since they are dictating where everybody is running. They will describe what they are looking for. The odds that such a place exists will be a test with DL between around -5 and 20, where 0 is obvious, 5 is easy, 10 is difficult, 20 is near impossible. Harder tests mean that the place is either very hard to find through some camouflage or are just very far away.

The fugitives can only describe one thing that they are looking for per round. If the place is not found, resume the chase with running and adjust the distances. The terrain feature is at a set distance tier from where the characters are, and is approached one tier per round. 

The terrain features can be: a bunch of objects to use as cover, buildings to hide in or barricade, a group of potential allies, a hole on the ground, a choke point, a river, a cloud of fog, a forest or anything else imaginable.

Some terrain features can be obstacles. Traversing an obstacle works just like running, but with a balance, climb or swim test to determine the ability to traverse. Again, use the lowest skill in the group to determine traversal ability.

\textbf{Hiding:} Hiding can be done at any distance greater than Combat in a Chase scene.

\textbf{Hiding at far distances:} If a hiding spot is entered at Far or Horizon distances, the chase is over. Whether the pursuers will be able to follow depends on interpretation of the scene in question. Maybe the fugitives have disappeared into the forest and need to be tracked, but maybe they are inside a house and cant leave without being seen. It all depends on the scene.

\textbf{Hiding at shooting distances:} The stealth test is done by the lowest stealth in the group, but splitting the group is possible. Stealth is done against the best of the opponent's detection. The hiding spot in question is relevant, since it can grant stealth bonuses. On a stealth miss or graze, the pursuers find the lowest stealth character and anyone that is with them. On a hit or crit, the pursuers follow some wrong track or give up.

\section{Stealth}

\skill{Stealth}{ None }{ Skill to remain unnoticed. It involves stepping lightly, moving with appropriate timing, keeping to the shadow, avoiding making noise. 

\textbf{Stealth value:} The player rolls the d10 and adds their stealth skill, that is the stealth value that is used against every detection skill value until a new roll is called. 

\textbf{Stealth events:} are a description given by the player of what exactly the character is trying to do with their stealth, which defines which senses the detectors are using and sets a bonus or penalty to the detection test. The usual values should stay between -15 and +15.

An event could be trying to hide in a specific spot, trying to pass through somewhere undetected or taking some kind of action without being identified or located. A stealth event should take into consideration what noise the footsteps make in a path (altered by how fast characters are moving), how difficult it is not to bump into anything, how hard it is to avoid line of sight, how hard it is to stop gear from making noise, how close to a detector the path lies, how good a hiding spot is and how long you will be searched for if staying there.

Each event adds a bonus or penalty to the stealth value relative to a detector, who uses the best sense available. On hits or crits, the character stays unnoticed. On grazes, active senses detect, but passive ones just draw the attention of active senses and possibly make characters go searching with the active ones. On misses, the character is detected. The stealth value is modified by events that change the situational bonus.

Call for a new stealth test when the scene changes into combat, chase or interpretation. One failed test means that stealth is over, so repeating multiple tests will eventually lead to failure. This is why only one test is called for a scene. The variation comes from the player's actions in stealth events.

\textbf{Event Chaos modifier:} Some stealth events can have highly unpredictable outcomes and others can be the opposite. Use safe tests or risky tests to represent this fact.

\textbf{Hidden roll:} Optionally, the GM could roll the dice for the player, add it to the character's stealth skill and keep the result, without revealing it. The player then describes what they are trying to do. For every major shift in event difficulty, the GM decides on a bonus and narrates what is happening. The player must interpret whether their dice result was good or not. They won't give up early if their stealth value came too low or just walk by everyone if the was too high.

\textbf{Group roll:} Group rolls are used in chases or interpreted scenes. Start with the worst stealth skill of any character in a group and add penalties for the amount of characters, gear carried, mounts, pets, cargo and vehicles as stealth event penalty.  

\textbf{Split Group:} In interpretation or chase scenes, the action depends on description and interpretation. If someone is detected, but the group is split, only the worst stealth group is detected.
}

\skill{Detection}{ Awareness}{ Detection is used to detect things in stealth or just for looking for things in general.

\textbf{Detection as defense:} Detection is used as a defense against stealth, meaning that it is used passively without rolling.

\textbf{Group detection:} a single Detection value should be used for a whole group of characters, using the best Detection in the group for everybody. The group can be separated as desired, but then finding the best detection may involve using different distance bands for each group.

\textbf{Senses:} Any detection event must target a specific sense. The basic usable senses are vision, hearing, smell and touch. In each detection attempt, a sense must be chosen to determine how the detection happens. Senses can be active or passive. 

\textbf{Active senses:} are used to keep the detector in aware, meaning that after detection, the target is tracked by the detector and can be reacted to until a new stealth event occurs.

\textbf{Passive senses} can detect position, but they can't track it. If the character has no active senses available, they need to detect every event from a single source to be able to react to it. The character is considered disoriented for a reaction if only passive senses are available to detect. Detecting through a passive sense can focus any available active senses on the target, detecting them actively.

\textbf{Vision:} is an active sense. Vision gets blocked by cover and requires light. The field of vision covers 180 degrees, which is defined in the grid as within the arc defined by the directions of 3 adjacent hexes around a character.

\textbf{Hearing:} is a passive sense and omnidirectional.

\textbf{Smell:} is a passive sense. It requires a smell trail to trigger the sense. Humans have -5 to this sense. Sniffing something from very close gives +5 to this test.

\textbf{Touch:} is an active sense. Anything that touches a character is immediately detected and the character becomes alarmed to it. Projectiles that hit the character have their direction identified, not necessarily their source. Using only touch as an active sense makes the character disoriented against another involved in the grappling, and oblivious against characters in melee range that have been detected only through touch.

\textbf{Distraction:} A character can be distracted by something, receiving -3 to Detection and cunning against everything other than the distraction. 
% A character that has used its action surge is automatically distracted.

\textbf{Looking for objects:} When trying to decide whether characters find an object, the object is found with a hit against the object's DL. If Detection is equal or higher than DL, they find it. 

\textbf{Investigation:} When a character is looking for something that requires specific knowledge background, they use the value of Detection as skill, but use knowledge as an enabler.

\textbf{Scanning:} Use detection to see how much you notice from a room. A good detection test reveals good paths for sneaking, hiding spots, dangers, and other information that will help players make good decisions when describing their stealth events.
}


\subsection{Guidelines on playing stealth}

These guidelines exist to provide a physical basis to draw the stealth bonuses. It is not meant to be used as a table, but to ground difficulty decisions and provide common sense to everyone at the table. In practice, the bonuses should be chosen quickly and intuitively. The guidelines are meant to sharpen instincts.

The difficulty to detect or perform stealth depend on which sense is the dominant one for that detection event. 

It is important to let players know how difficult a stealth event is expected to be. It is not necessary to give them the exact number for the stealth test, but at least an idea of what their character should expect from attempting a test.

If a player gives a description of something that absolutely should work, just let it work without a test. But also ask yourself if their character would come up with that. If the answer is more likely to be no, ask for a detection check, or another skill that may fit, like cunning, exploration or insight, to see if the character really would get that idea. If the character does not pass the test, they did not have that idea or they get a dumbed down version of it. 

The factors that determine the difficulty to detect or hide are: distance, noise, signal strength and time.

\textbf{Distance bands:} Detection is penalized by distance. It is convenient to design distance bands, as shown in the list. For each band over the first, add -3 penalty to vision, -5 to hearing and -5 to smell.

\begin{proplist}
    \item immediate (within 10m), bonus instead of penalty
    \item combat (within 50m), no penalty
    \item shooting(within 200 m)
    \item far (within 500 m)
    \item horizon
\end{proplist}

\textbf{Noise:} is the amount of conflicting signal in a given situation. For hearing, its the loudness and rhythm of other sounds. For vision, it is the amount of camouflage, as well as the amount of cover and the presence of very bright nearby objects. For smell, it is the variety and intensity of other smells. 

\textbf{Signal Strength:} This is what the character is trying to minimize with their stealth skill. It is based on the character's gear and behavior. Items that help with stealth should influence the signal strength component. The item should provide a description of how it helps.

Here's a list of possible modifiers. Many modifiers can be combined.

\textbf{Hiding in stealth:}
\begin{proplist}
    \item \textbf{Being in plain sight:} auto detect unless camouflaged or at far distance and beyond.
    \item \textbf{Camouflage in plain sight:} +3 to Stealth vs vision. Varies with camouglage quality.
    \item \textbf{Partial cover:} +3 to Stealth vs. vision.
    \item \textbf{Complete cover:} Automatically undetected to vision.
\end{proplist}

\textbf{Moving in stealth:} 
\begin{proplist}
    \item \textbf{Movement in the grass:} 0 vs. vision, -2 to Stealth vs. hearing.
    \item \textbf{Movement on marble floor with standard boots:} -2 to Stealth vs. hearing.
    \item \textbf{Movement on marble floor with muffled boots:} +0 to Stealth vs. hearing.
    \item \textbf{Wearing noisy gear or carrying jingly objects:} -2 to -4 to Stealth vs. hearing. Metal armor is around -2.
    \item \textbf{Wearing metallic armor:} -2 to Stealth vs. hearing.
    \item \textbf{Carrying burden:} -2/-4/-10 to Stealth vs. hearing for medium, large and cargo burdens.
    \item \textbf{Moving on squeaking floorboards:} -3 to -5 to Stealth vs. hearing.
    \item \textbf{Careful movement and crawling:} 0 to hearing and vision (Base Skill).
    \item \textbf{Basic movement:} -2 to Stealth vs. hearing and vision.
    \item \textbf{Running / Combat Sprint:} -5 to Stealth vs. hearing and vision.
\end{proplist}
    
\textbf{Attacking from stealth:}
\begin{proplist}
    \item \textbf{Attacking in melee:} -3 to Stealth vs. hearing. Causes immediate automatic detection if within a detector's field of vision.
    \item \textbf{Hitting in melee:} detected by touch. Detected by vision a moment later.
    \item \textbf{Shooting a bow:} -2 to Stealth vs. hearing and vision.
    \item \textbf{Hitting at distance:} -2 vs Vision. Direction is detected either way, so stealth cannot hit or crit.
\end{proplist}

\magic spellcasting is naturally undetectable without special detection. Only the spell's effects are detectable.

\textbf{Luck:} The longer one stays somewhere and someone looks for them, the easier it is to detect them. The greater the number of people searching and how much they are moving also factors into the odds of detection. The more distracted everyone is, the lower the chances of detection. Use luck to represent the character's lingering somewhere and giving luck a chance to ruin their stealth.

\subsection{Stealth in combat}

Stealth should not be free form in combat. The rules in combat situations are stricter to allow tactical planning.

\textbf{Stealth test:} The stealth value is rolled in the beginning of the round for each character in stealth individually. This is only done for characters that started their round undetectable to any opponent active senses. Enemies cannot target the character unless their detection is higher than the stealth value (a miss in stealth), and cannot react to them if their detection is not 4 points lower, at most (a graze in stealth). The stealth value of any character in combat is public to everyone and must be checked against whenever someone tries to act or react against a stealthy character.

\textbf{Stealth penalties:} The stealthy character's actions change their stealth value. Their stealth value situational bonus is set as the sum of all penalties they received until a new stealth test is made. The situational bonus in combat are:

\begin{proplist}
    \item \textbf{Tier 1 - Careful movement and crawling:} 0, \textbf{Tier 2 - Normal movement:} -2, \textbf{Tier 3 - Running and jumping:} -4.
    \item \textbf{Ranged attacks:} -2. \textbf{Melee attack:} -4.
\end{proplist}

\textbf{Staying in stealth:} The character must begin their turn undetectable through any active sense in order to be considered in stealth. They must go back to being undetectable to any enemy active senses and end their turn to continue in stealth. Failing to do so causes automatic detection by everyone (their stealth value disappears). Trying to take a second turn immediately after finishing their last one takes -5 penalty to cunning.

\textbf{Attacking from stealth:} Being attacked automatically causes detection through touch after the attack, which means that a second attack can always be reacted to. Whoever is struck always knows where the attack came from, but cannot target the spot directly unless they detect the attacker through an active sense or has high enough detection. 

\textbf{Searching actively:} A character can spend 3 AP to pay better attention to their surroundings. This increases their detection by +3 for the round. 

\textbf{Staying undetectable:} A character must use cover, fog, darkness or camouflage to stay undetectable. 

\textbf{Camouflage:} is necessary to stay stealthy in plain sight, and must be matched to the current environment. A camouflage item comes with a penalty to stealth that is added when the character is in plain sight. 

\textbf{Floor noise:} The terrain may have a bonus or penalty to stealth based on what the floor is made of, which modify the penalties from movement. Some floors make noise when stepped on, some halls are echoey, some floors have leaves and branches that can make noise when stepped on. This must be drawn in the map and the GM must inform the players of the numerical values associated with the terrain. For example: Marble floor makes movement penalties to stealth increase in 1 per movement tier.

\textbf{Gear noise:} Some gear produces noise when moved or worn, which will be indicated in the gear description. When it does, it increases the penalties from movement by +1 per tier and also the penalties from attacking by +1.

\section{Morale}

Morale is the measure of the characters' stress tolerance. Stressful situations, such as combat and exploration in terrifying places, trigger morale tests. The test is conducted using the will skill.

\textbf{Combat morale test:} At the beginning of each round, all characters make a will test against the DL of the terror of the moment. A graze leaves the character frightened, and a miss activates the character's limit stress action. The limit stress action only triggers if there is an enemy within 30m.

\textbf{Combat morale DL = aggravating factors}

\begin{proplist}
    \item \textbf{Sensory deprivation (being oblivious):} +3
    \item \textbf{Reduced sensory quality (being disoriented):} +2
    \item \textbf{Injured:} +2 per injury penalty
    \item \textbf{Burning or suffocating:} +5
    \item \textbf{Out of STA:} +2
    \item \textbf{Monstrous presence:} variable. Depends on presence.
\end{proplist}

\subsection{Social actions}

Every social action costs 4 AP +1 STA and counts as a Rest action. They only work on creatures capable of communication.

\textbf{Intimidation:} A character can make an intimidation attempt in combat. This increases the DL of the enemies next morale test of the $1 + character's charisma/2$. This can be done during a combat surge.

\textbf{Taunt:} In combat, a taunt modifies the next morale test to make enemies enraged instead of afraid. This increases the DL of the enemies next morale test of the $1 + character's charisma/2$.

Multiple intimidations and taunts don't stack. Only the highest one counts.

\textbf{Rally:} Say something encouraging in combat. This adds decreases the DL of allies next morale test of the 1 + character's charisma/2. This can be done during a reaction surge. Multiple intimidations don't stack. Only the highest one counts.

\textbf{Surrender:} An attempt to surrender. It is a persuasion test against the highest insight in the enemy team. On a hit or crit, nobody in the enemy team is allowed to target this character. Being perceived holding a weapon, attacking anyone, trying to hide, running away or doing anything to help their own team cancels the effect and increases the difficulty of surrendering by 5 for any subsequent attempts of anyone in the whole team.

\textbf{Negotiation:} This is a persuasion test against the highest between insight or cunning of an enemy. It is a social test and adds IB, just like as explained in the Story chapter. Only intelligent creatures capable of communication can accept this. The point of persuasion is to propose a truce for further conversation or demand capitulation. 

\textbf{Feign Death:} The character becomes immobile and prone. Anyone trying to target them has to hit or crit on an insight or cunning test against deception + IB, or be forced to do something else. The IB equals -5 + injury penalty (injury makes IB better). Feign death only works once per combat. If the character is attacked, they can break out of immobility and react normally.

\textbf{Call help:} Shouts for help, triggering a detection test for anyone outside the combat area. The DL is 0 - d6, or a negative d6. If help hits on the detection check, they move in to help within 2 turns if they were in shooting range.

%----------------

\end{multicols}